% ============================================================
%  APLICACION MOVIL EN UNITY
% ============================================================

\section{Aplicación Móvil en Unity}

La aplicación Android constituye el cliente interactivo del sistema. Su función es
recibir los datos del estudiante, comunicarse con el servidor, transformar las
soluciones numéricas en una representación tridimensional y administrar la interacción
táctil y sensorial. La resolución por diferencias finitas se ejecuta exclusivamente en
Go; el cliente no contiene una segunda implementación del método numérico. Esta
separación evita resultados divergentes, reduce la carga del teléfono y permite
actualizar los solvers sin recompilar el APK.

El teléfono Android del estudiante cumple la función de IHM del proyecto. Se eligió
porque constituye un recurso disponible
para la totalidad de los alumnos de la cátedra, lo que permite desarrollar la actividad
sin exigir la adquisición ni la disponibilidad de equipamiento adicional. La pantalla
táctil recibe la identificación, los parámetros y la selección de nodos, mientras que
la unidad de medición inercial (IMU) aporta la orientación utilizada para explorar la
escena tridimensional. Esta decisión elimina la necesidad de un periférico externo y
facilita el uso de la plataforma en la cátedra de Matemática Avanzada mediante un
dispositivo de uso cotidiano. De manera complementaria, la aplicación conserva un
modo estereoscópico para Google Cardboard. Este modo no altera el requisito principal:
el teléfono por sí solo continúa siendo suficiente para realizar la actividad.

\subsection{Entorno de desarrollo y organización por escenas}

El proyecto se desarrolló con \textbf{Unity 6.4}, versión
\texttt{6000.4.2f1}. Utiliza Universal Render Pipeline \texttt{17.4.0}, Input System
\texttt{1.19.0}, uGUI \texttt{2.0.0}, TextMeshPro y Google Cardboard XR Plugin
\cite{unityManual,cardboardQuickstart}. La
dependencia de Cardboard permite ofrecer dos formas de visualización seleccionables
durante la ejecución: pantalla plana monoscópica y pantalla estereoscópica para visor.

El ejecutable incluye tres escenas:

\begin{enumerate}
    \item \textbf{MenuScene}: identificación, descubrimiento y registro.
    \item \textbf{HelloCardboard}: ecuación de onda bidimensional.
    \item \textbf{HeatPlate}: ecuación de calor bidimensional.
\end{enumerate}

La escena inicial es plana y táctil. Una vez confirmado el registro se carga la onda.
Desde cada simulación existe un botón para cambiar de modelo, y la sincronización puede
producir el mismo cambio automáticamente.

Los componentes principales se resumen en la Tabla~\ref{tab:componentes-unity}.

\begin{table}[H]
\centering
\caption{Componentes principales de la aplicación Unity.}
\label{tab:componentes-unity}
{\small
\begin{tabular}{ll}
\toprule
\parbox[t]{4.3cm}{\textbf{Componente}} & \parbox[t]{9.2cm}{\textbf{Responsabilidad}} \\
\midrule
\parbox[t]{4.3cm}{\texttt{MenuController}} & \parbox[t]{9.2cm}{Descubrimiento UDP, registro HTTP y acceso.} \\
\parbox[t]{4.3cm}{\texttt{ResponsiveMenuUI}} & \parbox[t]{9.2cm}{Construcción adaptable de la pantalla inicial.} \\
\parbox[t]{4.3cm}{\texttt{CardboardStartup}} & \parbox[t]{9.2cm}{Arranque condicional de la cámara plana o del subsistema XR.} \\
\parbox[t]{4.3cm}{\texttt{DisplayModeController}} & \parbox[t]{9.2cm}{Selección persistente entre pantalla plana y Google Cardboard.} \\
\parbox[t]{4.3cm}{\texttt{MobileGyroscopeCamera}} & \parbox[t]{9.2cm}{Lectura, calibración y suavizado de sensores.} \\
\parbox[t]{4.3cm}{\texttt{Membrana2D}} & \parbox[t]{9.2cm}{Petición y reproducción de la onda.} \\
\parbox[t]{4.3cm}{\texttt{WaveParameterPanel}} & \parbox[t]{9.2cm}{Parámetros, altura y edición de nodos.} \\
\parbox[t]{4.3cm}{\texttt{HeatPlate2D}} & \parbox[t]{9.2cm}{Petición y reproducción del calor.} \\
\parbox[t]{4.3cm}{\texttt{HeatParameterPanel}} & \parbox[t]{9.2cm}{Parámetros térmicos y navegación.} \\
\parbox[t]{4.3cm}{\raggedright\texttt{SimulationSync}\allowbreak\texttt{Coordinator}} & \parbox[t]{9.2cm}{Seguimiento del estado docente.} \\
\bottomrule
\end{tabular}}
\end{table}

Los paneles se construyen durante la ejecución. Esto concentra estilos, anclajes y
eventos en componentes C\#, y reduce la dependencia de jerarquías serializadas.

\subsection{Modos de visualización y sensores}

La aplicación inicia el menú como una interfaz plana y, dentro de las escenas de
simulación, permite alternar entre dos modos. El modo predeterminado presenta una sola
imagen a pantalla completa y utiliza la IMU del teléfono, por lo cual no requiere un
visor. El modo VR inicializa Google Cardboard XR, habilita el renderizado estereoscópico
de URP y genera una imagen corregida para cada ojo. La selección se conserva al cambiar
entre las escenas de onda y calor.

\begin{lstlisting}[style=csharp, caption={Selección del modo de visualización.}]
XRManagerSettings manager = XRGeneralSettings.Instance != null
    ? XRGeneralSettings.Instance.Manager
    : null;

if (DisplayModeController.IsVrModeRequested)
{
    yield return StartCardboard(manager, mainCamera, displayMode);
}
else
{
    yield return StartFlatMode(manager, mainCamera, displayMode);
}
\end{lstlisting}

En Android se prioriza \texttt{GameRotationVector}, que estima la orientación sin
depender del norte magnético. Si no está disponible se usa \texttt{RotationVector};
como último recurso se integran las velocidades angulares del giroscopio
\cite{androidSensors}.

Al ingresar se guarda la rotación inicial y se espera un intervalo breve antes de
fijar la referencia. La rotación relativa es

\begin{equation}
q_r=q_0^{-1}q,
\end{equation}

donde $q$ es la actitud actual y $q_0$ la referencia. El orden expresa el movimiento
en el sistema local inicial y evita intercambiar cabeceo y giro lateral en horizontal.
Se conservan estos dos movimientos y se fija el alabeo en cero:

\begin{equation}
q_c=q_i\,q(\theta_x,\theta_y,0).
\end{equation}

La interpolación esférica utiliza

\begin{equation}
\lambda=1-e^{-s\Delta t},
\end{equation}

Aquí $q_i$ representa la orientación inicial de la cámara en la escena, antes de aplicar
el movimiento relativo medido por el teléfono. El orden de multiplicación conserva esa
orientación como referencia y aplica luego los ángulos relativos de cabeceo y giro
lateral. En la expresión de suavizado, $s$ controla la respuesta del filtro. Así se
mantiene la orientación original de la escena
sin perder una respuesta natural al movimiento del teléfono.

El seguimiento implementado es rotacional, con tres grados de libertad (3DoF): permite
observar la escena al modificar la orientación, pero no determina la posición del
teléfono en el espacio como lo haría un sistema de seis grados de libertad (6DoF). El
campo de visión y la corrección óptica dependen del perfil del visor configurado por
Cardboard, por lo que la aplicación no fija un valor angular único.

En modo VR, el seguimiento de cabeza es provisto por Cardboard mediante el componente
de seguimiento de pose de Unity. Los paneles táctiles planos se ocultan para no superponerse
a las imágenes estereoscópicas y la cruz nativa de Cardboard permite regresar al modo
plano. Esta variante deja preparada la visualización para quienes dispongan de un visor.
La conexión de un mando externo y el mapeo de sus botones para recorrer paneles,
seleccionar nodos y modificar parámetros se consideran una implementación futura; por
lo tanto, el informe no los presenta como una capacidad concluida de la versión actual.

\subsection{Menú de inicio adaptable}

La pantalla de ingreso funciona en posición horizontal y se adapta a diferentes
resoluciones, relaciones de aspecto, esquinas redondeadas y recortes de cámara. La
columna izquierda presenta la identidad del simulador y una onda animada; la derecha
contiene una tarjeta con nombre, legajo, estado del servidor y conexión.

Fondo, tarjeta, bordes y onda se generan con componentes gráficos de Unity. No dependen
de una captura fija y conservan nitidez. El \texttt{CanvasScaler} utiliza una referencia
horizontal de $1920\times1080$ y combina ancho y altura:

\begin{lstlisting}[style=csharp, caption={Escalado adaptable del menú.}]
scaler.uiScaleMode = CanvasScaler.ScaleMode.ScaleWithScreenSize;
scaler.referenceResolution = new Vector2(1920f, 1080f);
scaler.screenMatchMode =
    CanvasScaler.ScreenMatchMode.MatchWidthOrHeight;
scaler.matchWidthOrHeight = 0.5f;
\end{lstlisting}

La zona segura de Android se transforma en anclajes normalizados:

\begin{equation}
\boldsymbol{a}_{\min}=\left(\frac{x_{\min}}{W},\frac{y_{\min}}{H}\right),
\qquad
\boldsymbol{a}_{\max}=\left(\frac{x_{\max}}{W},\frac{y_{\max}}{H}\right).
\end{equation}

Para pantallas panorámicas se mantienen dos columnas equilibradas. En formatos
cercanos a $4:3$, la identidad visual se estrecha y la tarjeta recibe mayor ancho.
Los textos usan autoajuste, mientras que campos y botones conservan alturas táctiles.

La IP manual permanece oculta durante el flujo normal. Un enlace la despliega cuando
la red bloquea el broadcast. Durante la conexión se deshabilitan los campos, el botón
cambia a \textit{CONECTANDO...} y un indicador muestra búsqueda, éxito o error.

\subsection{Descubrimiento, registro y comunicación}

Sin dirección manual, el cliente transmite
\texttt{RVPROYECTO\_DISCOVER\_V1} por broadcast UDP al puerto \texttt{47777}. La
respuesta identifica servicio y puerto HTTP; la IP se obtiene del origen del datagrama.
La búsqueda espera $4{,}5$ segundos y requiere una red local común.

Luego se envía \texttt{POST /register} con nombre y legajo. Dirección y datos se
guardan en \texttt{PlayerPrefs}. La onda se carga solamente después de una respuesta
satisfactoria.

El descubrimiento y las peticiones HTTP fueron diseñados para una red local de aula.
El tráfico no se cifra y \texttt{PlayerPrefs} se utiliza como almacenamiento de
conveniencia, no como un mecanismo destinado a proteger información sensible. En
consecuencia, la operación prevista requiere una red Wi-Fi conocida y controlada.

Las simulaciones se solicitan con \texttt{UnityWebRequest} y JSON. El manifiesto
Android declara \texttt{INTERNET} y \texttt{ACCESS\_NETWORK\_STATE}; además habilita
HTTP local sin cifrar. El cliente conserva únicamente el cambio más reciente y espera
$5{,}2$ segundos entre cálculos. Ante HTTP 429 reintenta automáticamente.

\subsection{Estructura de datos y reproducción}

Onda y calor reciben ancho, alto, número de cuadros y un arreglo lineal. Unity verifica

\begin{equation}
N_{valores}=N_xN_yN_t.
\end{equation}

El cuadro $k$ y nodo $i$ se localizan mediante

\begin{equation}
j=k(N_xN_y)+i.
\end{equation}

Si las dimensiones no coinciden, la respuesta se rechaza. La secuencia válida se
reproduce a 30 cuadros por segundo con un acumulador basado en
\texttt{Time.\allowbreak deltaTime}.

Cada cuadro actualiza dos texturas: \texttt{\_MainTex} contiene el color y
\texttt{\_Displacement} la elevación. Separarlas impide que el color científico altere
la geometría.

\begin{lstlisting}[style=csharp, caption={Validación de una grilla recibida.}]
long expected = (long)result.width
    * result.height * result.frameCount;

if (result.width <= 0 || result.height <= 0 ||
    result.frameCount <= 0 ||
    expected != result.values.LongLength)
{
    return false;
}
\end{lstlisting}

\subsection{Escena de la ecuación de onda}

\texttt{Membrana2D} usa \texttt{POST /solve/wave2d}. El panel modifica $L$, $T$, $c$
y condición gaussiana, sinusoidal, triangular o modal. Los modos usan índices $a$ y
$b$ entre 1 y 5. También se habilitan amortiguamiento y excitación sinusoidal. La
frecuencia natural mostrada es

\begin{equation}
f_{ab}=\frac{c}{2L}\sqrt{a^2+b^2},
\end{equation}

y la excitación enviada es $f_e=m_f f_{ab}$. Esto permite explorar la resonancia.

La escala cromática usa azul para desplazamientos negativos, verde cerca del equilibrio
y rojo para positivos. Se mezcla una cuadrícula clara y se dibujan ambas caras para
observar la placa desde arriba o abajo.

Cada vértice posee un punto blanco implementado con partículas. El punto muestrea la
misma textura en la misma coordenada UV que el material, por lo cual permanece adherido
a la superficie durante la evolución.

\subsubsection{Condición inicial personalizada}

El usuario puede pausar y editar nodos mediante el retículo central y el giro del
teléfono. La elevación se transforma en desplazamiento y se extiende a vecinos dentro
de un radio de pincel. Al confirmar se envían \texttt{initialWidth},
\texttt{initialHeight} e \texttt{initialValues}. El servidor interpola la grilla y
calcula la evolución. Ante una falla, la forma se conserva para reintentar.

El panel y el control de altura pueden ocultarse. Los parámetros se envían después de
un breve período sin escritura para evitar una petición por carácter.

\subsection{Escena de la ecuación de calor}

\texttt{HeatPlate2D} usa \texttt{POST /solve/heat2d}. El usuario controla $L$, $T$,
difusividad $\alpha$ y condición inicial gaussiana, de bordes o de placa.

Las temperaturas se normalizan en $[0,1]$. La escala recorre oscuro, amarillo, naranja
y rojo. También se agrega una grilla y elevación moderada:

\begin{equation}
h=\operatorname{sat}\left(0{,}5+\eta T\right).
\end{equation}

El panel permite cambiar la altura y regresar a la onda.

\subsection{Sincronización docente}

\texttt{SimulationSyncCoordinator} persiste mediante \texttt{DontDestroyOnLoad} y
consulta \texttt{GET /sync/status} cada $1{,}5$ segundos. La respuesta indica actividad,
ecuación, parámetros y versión incremental.

Si la ecuación publicada no coincide con la escena actual, Unity cambia de escena. Si
coincide pero la versión cambió, solicita la nueva solución. El sondeo no ejecuta el
solver por sí mismo y las versiones evitan aplicar repetidamente una orden.

\begin{lstlisting}[style=csharp, caption={Cambio de escena sincronizado.}]
string targetScene = status.ecuacion == "calor2d"
    ? HeatScene
    : WaveScene;

if (SceneManager.GetActiveScene().name != targetScene)
{
    lastAppliedVersion = status.version;
    SceneManager.LoadScene(targetScene);
    yield break;
}
\end{lstlisting}

Una interrupción no detiene la animación recibida: se omite esa consulta y se reintenta
en el siguiente intervalo. Así, el profesor propaga un modelo común sin conexiones
persistentes.
